% !TeX root = ../thuthesis-example.tex
\chapter{绪论}
\section{研究背景及意义}
数据显示，2022 年全球无人机产业的市场规模已超过 300 亿美元，并预计将在 2030 年扩展至 550 亿美元~\cite{nonami2025overview}。无人机已广泛应用于遥感、摄影测量、包裹投递等多个领域~\cite{benarbia2021literature}。无人机具有极高的机动性，能够到达不受地形限制的区域；同时，相较于地面移动机器人，无人机通常具备更高的运动速度。近年来，无人机凭借其机动性强、部署灵活以及对地形依赖程度低等优势，在航拍测绘、环境监测、农业植保、基础设施巡检等领域得到广泛应用~\cite{COLOMINA201479,ahirwar2019application}。随着传感器与控制技术的不断发展，无人机逐渐从“短时任务执行工具”向“长期自主作业平台”演进。

然而，受限于现有电池技术的能量密度水平，无人机的续航能力始终是制约其应用拓展的核心瓶颈之一~\cite{kim2018drone}。在实际任务中，大量应用场景（如定点监测、目标观察与数据中继）并不要求无人机持续进行空间位移，而是需要其长时间停留在目标区域上方。此类任务通常依赖悬停方式完成，而悬停过程中无人机需要持续输出推力以抵消自身重力。尽管飞行器在空间中保持静止，却仍然消耗大量能量，对于系统整体静态功率消耗非常高，导致能源利用效率显著偏低~\cite{beigi2022overview,cabuk2024modeling}。

\begin{figure}[H]
  \centering
  \includegraphics[width = \textwidth]{kinds_of_UAV.png}
  \caption{典型的无人机系统}
\end{figure}

除续航问题外，无人机在安全性与任务鲁棒性方面同样面临挑战。悬停状态下，无人机的推力裕度有限，一旦受到风扰、气流不稳定或发生轻微碰撞，系统往往难以及时恢复稳定飞行，存在较高的失稳风险~\cite{schmidt2016stability}。此外，在复杂环境中执行任务时，无人机需要在机动性与飞行平稳性之间进行权衡，而单纯依赖飞行控制策略往往难以同时满足两方面需求。

针对上述问题，可栖息无人机系统逐渐成为研究热点。通过主动栖息于环境中的结构表面（如树枝、墙面、天花板或其他人工或自然物体）~\cite{hauf2023learning,lan2024aerial,li2025treecreeper}，无人机可在无需持续悬停的情况下完成定点任务，从而显著降低能耗、延长任务执行时间，并提升系统在复杂环境下的稳定性与安全性。此外，栖息能力还拓展了无人机的起降与作业方式，使其不再依赖传统的水平地面条件，不再需要约束对地平动速度为零，为复杂环境下的应用提供了新的可能性~\cite{wuest2024agile}。

目前，大部分商业无人机生产者为解决续航问题，提出的解决方案主要为调整电池与自重的比重，或是采用功率密度更高的电池，部分特殊场景下也会采用连接线缆直接供电的方案（如消防无人机）~\cite{mohsan2022towards}。这些方案对提升续航能力的效果有限，或是极大的减小了无人机的灵活程度。

本文调研了目前部分市面在售主流产品（如深圳市大疆创新科技有限公司的部分在售商品等），总结如表\ref{tab:drone_comparison}所示。

\begin{table}[htbp]
\centering
\caption{典型消费级与轻量级无人机参数对比}
\label{tab:drone_comparison}
\begin{tabular}{lcc}
\toprule
无人机型号 & 续航时间（min） & 自重（g） \\
\midrule
DJI Mini 5 Pro   & 52 & 250 \\
DJI Mini 4 Pro   & 45 & 248 \\
DJI Mavic 3 Pro  & 43 & 958 \\
DJI Mavic 3      & 46 & 895 \\
DJI Mini 4       & 34 & 249 \\
DJI Air 3S       & 45 & 724 \\
DJI Neo          & 18 & 135 \\
Crazyflie 2.1    & 7  & 27 \\
Autel EVO Lite   & 40 & 866 \\
Parrot ANAFI     & 25 & 320 \\
\bottomrule
\end{tabular}
\end{table}

因此，围绕柔性附着结构与无人机系统的协同设计，开展微型无人机的自主栖息与脱附控制研究，对于提升无人机续航能力、增强系统安全性并拓展其应用场景具有重要的工程意义和研究价值。

\section{国内外研究现状}
现有关于无人机栖息技术的研究主要集中在柔性附着结构设计、无人机系统集成以及栖息控制策略等方面。

\subsection{栖息结构的设计}
从更广义的系统行为角度来看，“栖息”并非人工飞行器所特有的工程概念，而是自然界中广泛存在的一类能量优化行为模式。对于鸟类而言，长时间持续飞行需要消耗大量代谢能量，因此其在执行觅食、观察等任务时，往往通过寻找稳定的支撑结构（如树枝、电线或岩壁）实现间歇性停驻，从而显著降低能量消耗并提升环境感知的稳定性。类似地，昆虫在飞行过程中也频繁利用墙面、叶片或其他微小结构进行短暂停留，以实现能量补偿、姿态调整或信息处理等功能~\cite{kovac2016learning,roderick2017touchdown}。这类“飞行—停驻”交替的运动模式，本质上体现了一种面向能量效率与任务持续性的自适应运动策略。

\begin{figure}[htbp]
  \centering
  \begin{subfigure}[b]{0.3\textwidth}
  \centering
  \includegraphics[width=\textwidth]{bird_perch01.jpg}
  \caption{鸟类栖息}
  \end{subfigure}
  \hfill
  \begin{subfigure}[b]{0.3\textwidth}
  \centering
  \includegraphics[width=\textwidth]{dragonfly_perch01.jpg}
  \caption{昆虫栖息}
  \end{subfigure}
  \hfill
  \begin{subfigure}[b]{0.3\textwidth}
  \centering
  \includegraphics[width=\textwidth]{bat_perch01.jpg}
  \caption{蝙蝠栖息}
  \end{subfigure}
  
  \caption{自然界动物栖息实例}
\end{figure}

从仿生学角度分析，上述自然栖息行为的核心优势在于其将高能耗的持续运动转化为“低能耗驻留 + 间歇性机动”的混合模式，从而在整体层面优化能量利用效率，同时增强系统在复杂环境中的稳定性与环境适应能力~\cite{zheng2023metamorphic}。特别是在鸟类的抓握式栖息与昆虫的表面附着行为中，生物体能够在极短时间内完成从动态飞行到稳定接触状态的切换，并在外界扰动存在的情况下维持可靠的附着稳定性，这一过程体现出高度鲁棒的动力学调节能力与高效的接触控制机制~\cite{hsiao2022mechanically}。

受上述自然现象启发，将“栖息”能力引入无人机系统设计中，为解决传统旋翼无人机在持续悬停任务中能耗过高、稳定性不足等问题提供了一种新的思路。通过赋予无人机在环境结构表面实现短时或长时附着的能力，使其能够在执行监测、观测或通信中继等任务时主动退出持续飞行状态，从而显著降低整体能耗并提升任务持续时间。这一类由自然界飞行生物启发而形成的工程方法，为无人机从“持续飞行平台”向“空地混合驻留平台”的演化提供了重要的理论与实践基础。

\subsubsection{现有栖息结构设计}
部分研究团队利用刚性机械结构，设计了具有夹持与抓握功能的仿生夹爪，为无人机实现静态栖息提供了可行路径。例如，Hang 等人设计了一种多功能折叠夹具，该夹具在常态下充当无人机的起落架，而在折叠后可作为夹爪使用，使无人机能够稳定栖附于台体边缘或直立杆件上\cite{hang2019perching}；Kirchgeorg 等人则提出一种外包络式结构，通过其特殊的几何形态使无人机能够嵌卡在错综复杂的树枝之间，实现被动式的固定\cite{kirchgeorg2021hedgehog}；Roderick 等人受生物启发，开发了一种仿鸟类爪部的抓握机构，依靠对细长杆件的包络与摩擦力实现无人机的稳固栖息\cite{roderick2021bird}。

\begin{figure}[htbp]
\centering
\includegraphics[width = 0.9\textwidth]{grab.png}
\caption{部分基于抓握的栖息无人机设计，(a)\cite{hang2019perching}, (b)\cite{kirchgeorg2021hedgehog}, (c)\cite{roderick2021bird}}
\end{figure}

上述设计均巧妙利用了无人机本体结构或附加机构，增强了无人机与复杂环境的物理交互及适应能力，从而显著降低了悬停时的能量消耗。以 Hang 的设计为例，其在直立杆上栖息时相比持续空中悬停可减少约 $69\%$ 的能耗。Kirchgeorg 与 Roderick 的方案则允许无人机在抓握稳固后完全关闭旋翼，使功耗接近为零。然而，这类抓握式栖息方式通常对目标物体的几何形状、尺寸及表面材质有较高要求，例如需要特定直径的杆件或具备足够摩擦力的表面，因此在泛化适应性与实际场景普适性方面仍存在局限。

此外，另一类研究则基于不同的黏附机理实现无人机栖息。Park 等人利用电黏附效应，通过高压静电场使无人机能够附着于光滑墙体等表面，实现了无需系留的“贴附”栖息\cite{park2020lightweight}；Zhang 等人设计了一款具有钩爪结构的顶端机构，使其能够钩挂于建筑物缝隙或粗糙表面下方，实现倒挂式悬停\cite{zhang2017spidermav}；Jim Ang 团队则采用真空吸盘作为吸附装置，利用大气压差产生强吸附力，使无人机可负载自重数倍的外部载荷并稳定吸附于平整表面\cite{ang2022drone}。

\begin{figure}[htbp]
\centering
\includegraphics[width = 0.9\textwidth]{sticker.png}
\caption{部分基于黏附的栖息无人机设计，(a)\cite{park2020lightweight}, (b)\cite{zhang2017spidermav}, (c)\cite{ang2022drone}}
\end{figure}

这类黏附式栖息策略拓宽了无人机可栖附的表面类型，尤其在人工建筑表面具有较好应用前景。然而，其性能往往受表面平整度、材质导电性、环境湿度等因素影响，且通常需要持续能源维持吸附力，系统复杂度与功耗之间的平衡仍是实际部署中需要进一步解决的问题。总体而言，现有栖息无人机设计在节能方面表现出显著优势，但在环境适应性、机构轻量化与控制系统可靠性等方面仍面临诸多挑战。

\begin{table}[htbp]
\centering
\caption{不同无人机栖息方式的性能对比}
\label{tab:perching_comparison}
\renewcommand{\arraystretch}{1.3}
\begin{tabular}{p{2.5cm} p{3.5cm} p{3.8cm} p{4.0cm}}
\toprule
\textbf{栖息方式} & \textbf{优点} & \textbf{缺点} & \textbf{适用场景} \\
\midrule

顶部夹爪栖息
& 可直接悬挂于杆状或边缘结构，栖息稳定性较高，对控制精度要求相对较低
& 对目标几何结构要求较高，仅适用于可抓握结构；机构复杂度较高，额外质量较大
& 栏杆、电缆、树枝等细长杆状目标 \\

下方夹爪栖息
& 可利用无人机下方空间完成抓取，结构相对简单，适用于平台降落后的固定
& 对接触位置与姿态精度要求较高，难以适应复杂表面；通常需额外接近与对准过程
& 平面边缘、平台支架、规则结构表面 \\

静电吸附栖息
& 结构轻量化程度高、功耗较低，对表面损伤小，可适用于较大面积附着
& 吸附力受表面材质、湿度及污染程度影响显著，对粗糙表面适应性有限
& 光滑墙面、玻璃、金属或绝缘平整表面 \\

气压吸盘栖息
& 吸附力较强，能够提供较高载荷能力与稳定性
& 需额外真空产生装置，功耗较高，对表面平整性与气密性要求较高
& 光滑且气密性良好的平整表面，如玻璃、瓷砖、金属面 \\

\bottomrule
\end{tabular}
\end{table}

\subsubsection{基于电流变液（VSERF）的栖息结构设计}

在面向无人机栖停的柔性附着与黏附机制研究中，研究者已探索了多种基于智能材料或仿生结构的解决方案，例如介电弹性体执行器、基于颗粒阻塞的正压可变刚度结构以及受自然界启发的仿生抓附机构等\cite{o2008review, crowley20223d}。相较于这些方法，基于电流变液的柔性附着结构展现出其独特优势：它具有毫秒级的快速响应能力、仅通过电场即可实现的简单控制方式，以及附着与脱附高度可逆的特性。这些优势使其特别适用于对重量、空间和能耗极为敏感的微型飞行器平台\cite{liu2012electrorheological}。

电流变液是一类典型的场致流变智能材料，一般由微米或纳米尺度的介电颗粒分散于绝缘载液中而成。其核心机理是电流变效应：当施加外加电场时，颗粒发生极化并在电场力作用下沿电场方向排列，瞬间形成链状乃至柱状的微观结构，从而导致整个流体的表观黏度急剧增大，实现从液态到类固态的可逆转变\cite{hao2001electrorheological}。这种状态转变的强度和速度可通过电场强度进行精确、快速的调节。

在无人机栖停的应用框架下，通常将ERF封装集成于足端或接触面模块内。在需要附着时施加电压，ERF迅速固化，在与目标表面接触的区域产生强大且可控的黏附力与摩擦力，使无人机能够稳定“吸附”于垂直、倾斜甚至倒置的多种表面。切断电压后，材料恢复液态，黏附力即刻消失，无人机得以近乎零阻力地脱附起飞。这种通过电信号直接控制“固-液”转变以实现即时黏附与分离的机制，为无人机在复杂环境中实现高能效、高机动性的动态栖停提供了全新的技术路径。

\begin{figure}[H]
\centering
\includegraphics[width = 0.5\textwidth]{ERF.png}
\caption{电流变液（ERF）机理解释。(a) ER粒子，(b) 非ER粒子\cite{hao1998mechanism}}
\end{figure}

为进一步提升栖停的持久性与可靠性，赵慧婵教授团队研发了一种新型电压可切换电流变流体（Voltage-Switchable ERF, VSERF）\cite{pan2022stretchable}。该材料在电场激活后，即使撤去外部电源，仍能长时间维持其固化状态与附着能力。实验表明，单次激活可实现超过25分钟的稳定附着，且材料在反复使用中累计工作时间超过64小时。这一突破性进展显著降低了维持栖停状态所需的持续电能，使得无人机仅凭机载有限能源即可实现超长时间静默驻留，极大拓展了其在持续监视、环境监测等长航时任务中的应用潜力。

\begin{figure}[H]
\centering
\includegraphics[width = \textwidth]{VSERF_Zhao.png}
\caption{基于VSERF的柔性附着机构及其在无人机栖停中的应用演示}
\end{figure}


\subsubsection{VSERF 的材料力学特性分析}
 
从材料力学角度深入分析 VSERF 的工作机理，有助于理解其在无人机栖息应用场景中的性能边界与工程约束。根据 Pan 等人的原始研究~\cite{pan2022stretchable}，VSERF 装置在电场强度从 $0\,\text{V/mm}$ 提升至 $750\,\text{V/mm}$ 时，可实现约 $14.8$ 倍的压缩刚度变化与 $3.5$ 倍的切向刚度变化。这一各向异性的刚度调节特性在工程应用中具有重要意义：压缩方向的高刚度可抵抗栖息时重力引起的法向变形，而切向方向的可调刚度则有助于在接触初期提供一定的顺应性，缓解接触冲击。
 
在流变学建模层面，经典电流变液的稳态剪切行为通常采用 Bingham 塑性模型加以描述：
\begin{align}
\tau = \tau_0(E) + \eta \dot{\gamma}, \quad \tau > \tau_0
\end{align}
其中 $\tau$ 为剪切应力，$\tau_0(E)$ 为依赖外加电场强度 $E$ 的动态屈服应力，$\eta$ 为塑性黏度，$\dot{\gamma}$ 为剪切速率~\cite{hao2001electrorheological}。然而，VSERF 在断电后保持固化状态的行为无法由标准 Bingham 模型完整描述，其内在机理涉及颗粒链的锁固与极化电荷的持久保持，目前认为与颗粒间的 Maxwell 应力和 Johnsen–Rahbeck 效应密切相关~\cite{hao1998mechanism}。对 VSERF 独特的"自锁"流变行为建立更为精确的本构模型，是支撑其工程量化设计的重要基础性工作。
 
在力学性能定量层面，VSERF 装置在激活状态下可承载自重约 $19$ 倍的外部载荷（约 $80\,\text{g}$），且在拉伸应变不超过 $20\%$、弯曲曲率不超过 $50\,\text{m}^{-1}$ 的条件下，其刚度调节能力不受影响~\cite{pan2022stretchable}。这一柔顺性使其能够在非激活状态下自适应地贴合珊瑚石、铅笔、三维打印曲面等复杂形态的目标表面，从而显著降低对接触姿态精度的要求，这对于欠驱动四旋翼平台在接近阶段难以精确控制末端位姿的固有缺陷，具有重要的工程补偿意义。
 
\subsubsection{VSERF 与其他智能材料的横向对比}
 
为进一步明确 VSERF 方案相较于同类智能材料与黏附机构的技术优势与应用边界，表~\ref{tab:material_comparison} 从响应速度、驱动能耗、表面适应性、断电保持能力以及系统集成复杂度等维度，对当前主流的可变刚度与主动黏附材料进行了系统性对比。

\vspace{10pt}
\begin{table}[htbp]
\centering
\caption{主流智能材料/黏附机制用于无人机栖息场景的性能对比}
\label{tab:material_comparison}
\renewcommand{\arraystretch}{1.3}
\begin{tabular}{p{\dimexpr0.17\textwidth-2\tabcolsep\relax}
p{\dimexpr0.16\textwidth-2\tabcolsep\relax}
p{\dimexpr0.17\textwidth-2\tabcolsep\relax}
p{\dimexpr0.16\textwidth-2\tabcolsep\relax}
p{\dimexpr0.17\textwidth-2\tabcolsep\relax}
p{\dimexpr0.17\textwidth-2\tabcolsep\relax}}
\toprule
\textbf{材料/机制} & \textbf{响应时间} & \textbf{维持功耗} & \textbf{表面适应性} & \textbf{断电保持} & \textbf{集成复杂度} \\
\midrule
VSERF（电流变液）
& 毫秒级（$<$10 ms）
& 极低（断电保持）
& 高（可贴合复杂曲面）
& \textbf{是}（$>$25 min）
& 低（仅需高压电源） \\
 
介电弹性体（DEA）
& 毫秒级
& 持续高压
& 中（需预拉伸）
& 否
& 中（机械预载荷复杂） \\
 
形状记忆合金（SMA）
& 秒至分钟级
& 加热功耗高
& 低（几何固定）
& 否（需持续加热）
& 高（热管理复杂） \\
 
颗粒阻塞（Jamming）
& 毫秒至秒级
& 持续气压
& 高（粒子填充自适应）
& 否
& 高（需气泵与管路） \\
 
电吸附（EA）
& 毫秒级
& 持续高压
& 中（对光滑面效果好）
& 否
& 低
\\
 
永磁吸附
& 无（被动）
& 零
& 低（仅限铁磁表面）
& 是（永久）
& 低（脱附需额外机构） \\
 
\bottomrule
\end{tabular}
\end{table}
 
由表~\ref{tab:material_comparison} 可以看出，VSERF 在"断电保持"这一关键指标上具有其他方案所不具备的独特优势。对于无人机的长时驻留任务而言，维持栖息状态所消耗的静态功率是影响整体续航的核心因素之一。传统电吸附方案需要在整个栖息周期内持续施加高压电场，其静态功耗虽然相较于旋翼悬停已大幅降低，但在数十分钟乃至数小时的长航时任务中，仍会对有限的机载电池造成不可忽视的负担。与之相比，VSERF 的"激活一次、持续附着"特性从根本上解耦了栖息时间与能量消耗之间的正比关系，使无人机能够以接近零待机功耗的方式实现超长时间静默驻留。
 
值得注意的是，颗粒阻塞方案（Particle Jamming）在表面适应性方面与 VSERF 相近，且在软体抓手领域已有成熟的工程应用~\cite{o2008review}；然而其需要持续的气压维持以防止颗粒松散，并且气泵与气路管道的引入不可避免地增加了系统质量与结构复杂度，在微型无人机平台上的集成难度较大。
 
 
综合来看，VSERF 在响应速度、表面适应性、断电保持与系统集成复杂度之间取得了较为均衡的性能折中，尤其适合对能量效率与轻量化要求严苛的微型无人机平台。
 
\subsubsection{VSERF 在无人机栖息系统中的工程挑战}
 
尽管 VSERF 展现出上述显著优势，其在工程应用层面仍面临若干需要重点关注与解决的挑战。
 
高压电源的轻量化设计。VSERF 的激活通常需要施加数百至数千伏的电场强度（典型工作电场约为 $750\,\text{V/mm}$，对应毫米级厚度器件的工作电压在数百伏量级）。在微型无人机平台上，机载高压电源模块的重量与体积直接影响有效载荷预算。如何在保证电压转换效率的前提下，实现高压电源电路的极致轻量化，是 VSERF 工程化应用的首要硬件挑战。已有研究在微型飞行器电吸附领域探索了轻量化高压发生器的设计方案~\cite{park2020lightweight}，其技术路线对于本系统具有重要参考价值。
 
附着力的定量表征与预测。VSERF 对复杂表面的附着力受目标表面粗糙度、材质、温湿度环境以及接触面积等多重因素影响，目前尚缺乏针对无人机典型栖息场景（如竖直墙面、倾斜天花板、粗糙树皮等）的系统性实验表征数据库。在实际部署前，需要对不同表面类型下的法向附着力与切向摩擦力进行充分的实验标定，并建立考虑接触姿态误差与表面特性变化的附着力预测模型，以支撑控制系统的自适应调节。
 
多次循环使用的稳定性与寿命。尽管 Pan 在文章中指出 VSERF 在重复使用条件下累计工作时间超过 $64$ 小时，但针对无人机频繁起降场景（每次栖息-脱附循环对材料造成的机械疲劳与化学老化）的长周期耐久性评估仍有待深入研究。此外，柔性封装材料在反复弯曲与拉伸过程中可能产生微裂纹，从而影响电场分布的均匀性与附着力的一致性，需要在材料制备与结构设计阶段进行更详细的设计与检验。
 
与飞行控制系统的时序协同。VSERF 的激活时序需要与无人机的接近轨迹、接触检测以及栖息后的旋翼停机指令高度同步。过早激活可能导致材料在接触前即进入固化状态，导致预应力不足，无法牢固的黏附；过晚激活则可能在接触冲击后尚未形成有效黏附力便已脱离目标表面。因此，设计可靠的接触检测机制（如基于足端力传感器或惯性测量单元的接触事件识别算法），并据此精确触发 VSERF 的激活指令，是实现高成功率自主栖息的关键控制工程问题。
 
综上所述，VSERF 凭借其独特的断电保持特性、毫秒级响应速度以及对复杂表面的良好适形能力，为微型无人机的低功耗长时驻留提供了极具前景的技术路径。与此同时，其在轻量化高压驱动、附着力定量表征、循环耐久性以及控制时序协同等方面所面临的工程挑战，也正是本研究需要重点攻克的核心科学与工程问题。

\subsection{无人机控制系统的设计}

在无人机系统设计方面，国内外研究普遍采用轻量化结构设计与高效动力系统配置相结合的方式，以提升无人机的机动性能与飞行稳定性。通过减小机体质量、优化结构布局以及选用高推重比电机与螺旋桨组合，四旋翼无人机能够在有限能量约束下实现较大的加速度和良好的姿态响应特性。同时，成熟的开源飞控系统（如 Betaflight、PX4 等）被广泛应用于微型无人机平台，为姿态稳定控制和轨迹跟踪提供了可靠基础\cite{betaflight,PX4}。这些飞控系统通常基于经典控制理论或其改进形式，通过对电机转速的高速闭环调节，实现对无人机姿态与位置的稳定控制。

目前，学界已提出多种基于视觉信息的无人机控制方法，其所利用的视觉信息形式包括仅包含二维图像特征的视觉信息以及融合深度或空间位姿信息的三维感知数据。针对仅依赖二维视觉信息的控制场景，Mademlis 等提出了一种基于 PID 的视觉控制方法，在缺乏 GPS 与三维空间信息的条件下，仅利用二维像素特征即可实现多视角目标跟拍~\cite{mademlis2024vision}。类似地，Passalis 等同样基于 PID 控制框架，在仅获取二维视觉信息的情况下实现了对高速运动自行车运动员的稳定跟踪~\cite{passalis2018efficient}。另一方面，在具备精确位姿感知能力的场景中，Galvane 等在结合动作捕捉系统的实验环境下（即无人机可实时获取高精度空间位姿信息）提出了一套面向电影场景拍摄的无人机控制方法，使无人机能够沿预定轨迹完成多视角自主拍摄，并对拍摄过程进行三维场景建模与重构~\cite{galvane2017automated}。此外，Paredes 等通过神经网络训练，使微型无人机（包括扑翼与旋翼平台）在仅依赖微型摄像头输入的条件下实现悬停、降落以及侧向机动等飞行动作~\cite{paredes2024fully}。上述研究表明，视觉感知与飞行控制的深度耦合能够有效支撑无人机目标跟踪、自主导航及复杂机动任务的实现，并为本文基于视觉伺服的微型无人机自主栖息控制方法提供了重要的技术参考与研究基础。

\begin{figure}[htbp]
\centering
\includegraphics[width = \textwidth]{vision_controller.png}
\caption{部分基于视觉的无人机控制器研究，(a)~\cite{mademlis2024vision},(b)~\cite{passalis2018efficient},(c)~\cite{galvane2017automated},(d)~\cite{paredes2024fully}}
\end{figure}

从运动机理上看，典型四旋翼无人机通过四个固定安装、转速可控的旋翼产生升力和反扭矩，其运动控制本质上依赖于旋翼推力的差分分配。通过调节各旋翼转速的组合，无人机可以独立控制机体的滚转角（roll）、俯仰角（pitch）和偏航角（yaw），并通过整体推力大小调节垂直方向的升降运动。然而，由于旋翼推力始终沿机体法向方向作用，四旋翼无人机在平移运动上必须依赖姿态倾斜来间接产生水平分力，因此其系统在严格意义上属于欠驱动系统，仅能在任一时刻直接控制四个自由度，而无法对六个自由度的位置与姿态进行完全独立调节。这一欠驱动特性对轨迹规划、精细接触控制以及复杂环境下的机动操作提出了更高要求。

从系统建模角度来看，四旋翼无人机可被视为一个刚体，其构型空间由位置与姿态共同描述，状态向量可表示为
\begin{align}
\mathbf{q} =
\begin{bmatrix}
\mathbf{p} \\
\boldsymbol{\eta}
\end{bmatrix}
=
\begin{bmatrix}
x & y & z & \phi & \theta & \psi
\end{bmatrix}^\mathsf{T},
\end{align}

其中 $\mathbf{p} = [x,\, y,\, z]^\mathsf{T}$ 表示质心在惯性坐标系中的位置，$\boldsymbol{\eta} = [\phi,\, \theta,\, \psi]^\mathsf{T}$ 分别表示滚转角、俯仰角和偏航角。因此，四旋翼无人机在三维空间中具有六个构型自由度。

四旋翼系统的控制输入来源于四个旋翼产生的推力，其控制输入向量通常定义为
\begin{align}
\mathbf{u} =
\begin{bmatrix}
u_1 & u_2 & u_3 & u_4
\end{bmatrix}^\mathsf{T},
\end{align}

其中 $u_1$ 表示总推力，$u_2, u_3, u_4$ 分别对应绕机体系 $x$、$y$、$z$ 轴的控制力矩。该输入与各电机转速之间满足非线性映射关系，但其维数始终为四。

由此可见，系统的可控输入维数满足
\begin{align}
\dim(\mathbf{u}) = 4 < \dim(\mathbf{q}) = 6,
\end{align}

表明四旋翼无人机属于典型的欠驱动系统。在该结构约束下，飞行器无法同时对六个自由度进行独立控制，其水平位置运动 $(x,y)$ 需通过姿态角 $(\phi,\theta)$ 的倾斜间接实现，从而在控制与规划层面引入了显著的耦合关系。

这一欠驱动特性在常规飞行任务中可通过分层控制结构有效处理，但在涉及精细接触与环境交互的栖息过程中，姿态调整、速度调节与接触力控制必须高度协调，使得控制问题呈现出更强的非线性和时序敏感性，也进一步增加了实现完全自主栖息的难度。

在实际应用中，尽管现有飞控系统已在姿态稳定和常规轨迹跟踪方面表现出较高成熟度\cite{xia2024cooperative}，但在涉及近距离环境交互的任务中仍存在明显局限。尤其是在无人机栖息过程中，飞行器需要在极小空间内完成精准的姿态对准、速度调节与接触时序控制，而此类操作对感知精度、控制带宽以及多模态传感融合能力均提出了较高要求。受限于机载算力、传感器配置和系统集成复杂度，现有微型无人机系统多依赖人工操控完成最终接近与附着过程，自主决策与自主控制能力仍然有限。这一现状在一定程度上制约了无人机在长期驻留与能量受限任务中的应用潜力。

\subsection{仿真平台选择}

为验证所提出的无人机自主栖息与脱附控制框架的有效性，仿真平台需要同时具备高精度动力学建模能力、稳定的接触建模性能以及足够的数值计算效率，以支撑飞行—接触—附着过程中的强非线性动态行为分析。因此，在系统实现之前，有必要对主流机器人仿真平台进行系统性比较与选型分析。本文主要考虑的仿真框架包括 Gazebo~\cite{koenig2004design}、Isaac Sim~\cite{zhou2024towards} 以及 MuJoCo~\cite{zakka2025mujoco,todorov2012mujoco} 三类代表性平台。

\begin{figure}[htbp]
  \centering
  \begin{subfigure}[b]{0.3\textwidth}
  \centering
  \includegraphics[width=\textwidth]{gazebo.png}
  \caption{Gazebo场景}
  \end{subfigure}
  \hfill
  \begin{subfigure}[b]{0.3\textwidth}
  \centering
  \includegraphics[width=\textwidth]{issac_sim.jpg}
  \caption{Isaac Sim场景}
  \end{subfigure}
  \hfill
  \begin{subfigure}[b]{0.3\textwidth}
  \centering
  \includegraphics[width=\textwidth]{mujoco.png}
  \caption{MuJoCo场景}
  \end{subfigure}
  
  \caption{仿真平台典型场景}
\end{figure}

Gazebo 作为ROS生态中应用最为广泛的开源仿真平台，其优势在于与ROS系统的高度集成性以及对多机器人系统的良好支持。然而，Gazebo 基于传统的刚体动力学引擎（如 ODE / Bullet），在处理高频接触动力学与刚性碰撞问题时容易出现数值稳定性不足、接触力抖动以及时间步长敏感等问题。此外，其接触模型通常依赖于简化的库仑摩擦近似，在涉及柔性附着与高频接触切换的场景中，难以提供足够精细的力学一致性，从而限制了其在精细操作与接触任务中的适用性。

NVIDIA Isaac Sim 基于 PhysX 引擎，并结合GPU加速与高保真渲染能力，在视觉仿真与数据生成任务中具有显著优势，尤其适用于学习型方法中的大规模数据采集与域随机化训练。然而，其系统复杂度较高，对计算资源（尤其是GPU显存与图形管线）的依赖较强，同时其物理建模层在高频控制回路下仍存在一定的数值耗散与接触不连续问题。此外，对于以控制精度与动力学一致性为核心目标的研究任务而言，其仿真开销与系统复杂度并不具有最优性价比。

相比之下，MuJoCo在连续动力学建模与接触稳定性方面表现出更为突出的优势。MuJoCo 基于解析形式的多刚体动力学建模方法，并采用隐式时间积分与优化式接触求解策略，从根本上提高了高频接触过程中的数值稳定性与可重复性。尤其在涉及接触约束切换、摩擦锥近似以及软接触建模时，MuJoCo 能够在较大时间步长条件下仍保持良好的能量一致性与动力学连续性，这对于无人机在栖息过程中涉及的“自由飞行—瞬时接触—附着保持”混合动力学系统尤为关键。

此外，MuJoCo 在计算效率方面具有显著优势，其紧凑的动力学表达形式与高效的求解器使其能够支持高频控制回路下的实时仿真，从而与无人机控制频率（通常为100–1000 Hz）保持一致。这一特性对于研究接触瞬态过程中的控制稳定性与冲击抑制尤为重要。同时，MuJoCo 对系统参数具有良好的可解释性与可调性，使得接触刚度、阻尼以及摩擦参数能够以连续方式进行建模与调节，从而便于进行系统性消融实验与动力学分析。

综合考虑仿真精度、数值稳定性、计算效率以及与控制任务的匹配程度，本文最终选用 MuJoCo 作为主要仿真平台，用于构建无人机—柔性附着机构的统一动力学模型，并对栖息与脱附过程中的接触动力学行为进行系统性分析与验证。

\subsection{当前研究存在的问题}

尽管现有研究在无人机栖息机构设计与基础控制方法方面已取得阶段性进展，但从完整的自主系统视角审视，相关技术在环境感知的自主性、动态过程的鲁棒性以及实际部署的可行性等方面，仍面临一系列深层次的挑战与瓶颈。

\textbf{（1）感知与环境理解层面。}对栖息目标的自主识别与高精度位姿估计能力尚显不足。多数现有方案基于栖息目标的几何、位置或姿态已知的强假设，或依赖于操作员人工选定并引导最终接近过程，难以适应真实场景中目标尺寸多变、空间朝向随机以及存在部分遮挡等复杂情况。受限于微型无人机平台搭载的传感器精度、视场角及机载计算资源，在完成远距离逼近后，于关键的近距离接触阶段往往难以持续、稳定地获取满足精准操控需求的高频高精度相对状态信息，这直接导致了最终接触动作的失败率居高不下~\cite{adams2000contact}。

针对上述问题，本研究拟采用视觉识别特定基准标记（ArUco/AprilTag）的方案，在厘米级定位精度约束下构建轻量化的相对位姿估计流程，并结合惯性测量单元（IMU）进行高频融合滤波，以在近距离接触阶段维持稳定的状态反馈。

\textbf{（2）动力学与控制层面。}无人机从自由飞行状态到与环境形成稳定接触的瞬态切换过程，存在显著的混合系统控制难题。栖息动作的执行使得系统动力学模型发生突变：从受空气动力学支配的欠驱动飞行模型，瞬间转变为受接触力、摩擦力和可能存在的黏附力共同约束的强耦合系统。在接触建立或脱离的瞬间，相互作用力的阶跃式变化极易诱发剧烈的姿态振荡与高度波动，甚至导致系统失稳。传统的飞行控制器通常基于连续、平滑的动力学假设设计，难以有效处理此类由离散接触事件触发的、高度非线性的瞬态动力学行为。

从混合系统理论的视角来看，该问题可形式化表述为具有离散模态切换的切换动力学系统：
\begin{align}
\dot{\mathbf{x}} = f_{\sigma(t)}(\mathbf{x}, \mathbf{u}), \quad \sigma(t) \in \{\text{飞行}, \text{接触}, \text{附着}\}
\end{align}
其中 $\sigma(t)$ 为离散模态切换信号，在接触事件发生时刻产生跳变，而各模态下的向量场 $f_\sigma$ 具有本质不同的结构。传统单一控制器难以同时保证所有模态下的稳定性与性能指标，亟需面向模态切换的统一控制框架。

针对上述问题，本研究拟采用基于有限状态机的分层控制架构，在自由飞行、最终接近与接触附着三个阶段分别设计相应的控制律，并通过接触检测信号触发平滑的状态切换，抑制接触冲击引起的瞬态扰动。

\textbf{（3）接触姿态约束与适应性方面。}现有多数栖息策略对接触瞬间的无人机姿态施加了严格限制。无论是机械抓握还是场致黏附，通常都要求无人机末端执行器沿特定方向（通常是表面法向）与目标实现精确对准。这一要求与四旋翼平台固有的欠驱动特性（无法在所有方向上独立产生加速度）相结合，构成了一个复杂的控制约束。在狭窄或结构复杂的空间中，同时满足精确的末端位姿对准、安全的接近速度规划以及实时的避障需求，极大地增加了运动规划与反馈控制的复杂度，限制了系统对不同空间构型目标的适应能力。

这一矛盾在水平天花板栖息场景中尤为突出：无人机需要在接近竖直倒置姿态时完成接触，而此时旋翼推力方向接近水平，位置控制能力急剧退化，系统极易失稳。本研究拟通过 VSERF 柔性附着结构的顺应性特性缓解姿态对准的精度需求，并在规划层面引入面向欠驱动约束的可行姿态廊道，在保证稳定性的前提下放宽末端位姿的精度要求。

\textbf{（4）机构与系统集成层面。}栖息功能模块的引入不可避免地带来了额外的质量、惯量与功耗负担，对平台本身的续航时间、机动性与负载能力产生负面影响。部分基于主动黏附或吸附的方案需要持续的能量输入以维持附着状态，从而部分抵消了栖息节能的初衷；而具有复杂传动机构的抓握装置则可能引入新的单点故障风险，影响系统整体可靠性。因此，如何在附着性能、机构复杂度、系统重量及能耗之间取得最优平衡，尚未形成普适性的系统化设计准则。

表~\ref{tab:integration_tradeoff} 从质量开销、静态功耗与机构复杂度三个维度，对典型栖息方案的系统集成代价进行了定性比较，以直观呈现当前设计空间中的主要权衡关系。

\begin{table}[htbp]
\centering
\caption{典型无人机栖息方案的系统集成代价定性比较}
\label{tab:integration_tradeoff}
\renewcommand{\arraystretch}{1.3}
\begin{tabular}{lcccc}
\toprule
\textbf{方案类型} & \textbf{额外质量} & \textbf{静态维持功耗} & \textbf{机构复杂度} & \textbf{典型代表} \\
\midrule
刚性夹爪抓握   & 高   & 低（可断电锁止）  & 高（传动机构多） & \cite{hang2019perching} \\
真空吸盘附着   & 中   & 高（需持续气泵）  & 中               & \cite{ang2022drone}      \\
电吸附附着     & 低   & 中（需持续高压）  & 低               & \cite{park2020lightweight} \\
VSERF 柔性附着 & 低   & \textbf{极低} & 低          & \cite{pan2022stretchable} \\
\bottomrule
\end{tabular}
\end{table}

针对上述问题，本研究选取 VSERF 柔性附着模块作为核心机构，并以轻量化一体化设计为原则，将高压驱动电路集成于无人机机架之内，在不显著增加系统质量的前提下实现可靠的栖息与脱附功能。

\textbf{（5）系统自主性层面。}当前绝大多数栖息演示仍严重依赖人工遥控或半自主的人机协同操作模式。在最终接近轨迹微调、接触姿态对齐以及触发时机选择等关键决策环节，仍需人类操作员介入。这不仅增加了操作员的认知负荷与培训成本，更从根本上限制了系统在动态、未知或非结构化环境中的自主行动能力与任务可扩展性。实现从环境感知、决策规划到精准执行的全流程自主栖息，是推动该技术从实验室演示走向实际工程应用必须跨越的核心障碍。

本研究的目标正是补全这一自主性缺口：通过视觉感知、分层控制与 VSERF 时序协同的统一设计，在不依赖人工干预的条件下实现完整的“感知—规划—执行—附着—脱附”闭环自主流程。

上述五个层面的问题并非相互独立，而是形成了复杂的耦合关系：感知精度不足直接放大了控制切换的难度，欠驱动约束进一步压缩了姿态容错空间，而机构集成的质量代价又反过来削弱了飞行机动性与续航裕量。因此，本研究的核心科学问题与工程挑战在于：如何在微型无人机平台严苛的感知、计算与能源约束下，协同发展鲁棒感知与状态估计方法、平滑处理动力学模式切换的自适应控制架构，以及轻量化低功耗智能栖息机构的一体化集成方案。将三者纳入统一的系统设计框架，能够从源头上突破现有技术瓶颈，进而实现无人机的自主、可靠与高效栖息。

\section{研究内容和研究工作计划}

\subsection{研究目标}
本研究面向微型无人机在复杂非结构化环境中执行长时间驻留与持续感知任务时所面临的续航受限、能量利用效率低以及系统鲁棒性不足等关键问题，围绕柔性附着结构与飞行控制系统的深度耦合与协同设计，旨在构建一套具备环境感知、自主决策与稳定执行能力的无人机自主栖息与脱附一体化控制框架，从而实现飞行—接触—驻留多模态任务条件下的统一系统级优化。具体研究内容与目标包括以下几个方面：

首先，面向复杂、非结构化环境中的广义栖息目标，研究适用于微型无人机平台的鲁棒感知与相对位姿估计方法，实现对不同几何形态、空间姿态及表面特性栖息目标的稳定识别与高精度定位，为后续接近控制与接触动作提供可靠的状态估计基础。在实现路径上，拟采用基于视觉标记的轻量化感知方案（如 ArUco、AprilTag 等结构化视觉编码），在保证计算效率与系统可部署性的同时，实现厘米级甚至亚厘米级的相对位姿估计精度，从而满足微型无人机在近距离接触阶段的控制精度需求\cite{joseph2024autonomous,huang2023indoor}。在此基础上，进一步可扩展至弱纹理环境与自然场景下的鲁棒视觉定位问题，为后续泛化应用奠定基础。

其次，针对无人机从自由飞行状态向接触附着状态转变过程中所面临的动力学模型突变与控制模式切换问题，建立能够统一描述飞行—接触—附着全过程的混合动力学建模与控制框架。在该框架下，将飞行阶段的连续欠驱动动力学系统与接触阶段的动力学系统进行统一表征，重点研究面向强非线性与不连续约束条件的鲁棒控制策略与切换逻辑设计，实现无人机在接触瞬间的姿态稳定保持、速度连续过渡以及接触冲击的有效抑制，从而降低由于动力学不一致性引起的系统振荡与失稳风险。

再次，围绕基于电流变液的柔性附着机构，研究其与无人机飞行控制系统之间的跨域协同机制，重点分析材料电场响应特性与系统级控制时序之间的耦合关系，充分发挥其快速响应、可逆相变以及低维持功耗等核心优势。在此基础上，实现附着力可调、接触界面柔顺化以及脱附过程可控化的统一设计，从而构建一种轻量化、低能耗且高可靠性的栖息与脱附执行单元，使其能够在严格资源约束的微型无人机平台上长期稳定运行。

最后，通过仿真验证与实物实验相结合的方式，对所提出的自主栖息系统进行系统性评估与性能验证，重点考察其在不同栖息表面类型、空间构型变化以及外部扰动条件下的稳定性、鲁棒性与能量效率表现，从而形成一套可复现、可扩展的无人机自主栖息技术方案。仿真方面，基于高保真物理引擎 MuJoCo 构建包含无人机本体动力学与柔性附着结构的统一仿真模型，对栖息与脱附过程中的接触力演化、状态切换过程及能量消耗特性进行定量分析\cite{todorov2012mujoco}。在仿真验证基础上，进一步开展实体系统实验，对系统在真实环境中的栖息成功率、动态稳定性以及单位任务能耗进行综合评估，从而验证所提出方法在工程应用层面的有效性与实用性。

\begin{figure}[H]
  \centering
  \includegraphics[width = 0.8\textwidth]{target.png}
  \caption{本项目的工作目标}
\end{figure}

\subsection{研究思路}

针对当前无人机栖息技术在自主性、鲁棒性与系统集成方面存在的突出问题，本文采用“结构—感知—控制—系统验证”协同推进的研究思路，从系统层面出发，逐步构建完整的无人机自主栖息解决方案。

在系统建模与仿真层面，对四旋翼无人机的飞行动力学特性及其欠驱动约束进行建模，并引入柔性附着机构的接触与黏附模型，构建能够覆盖自由飞行、近距离接近以及接触附着阶段的统一仿真环境。通过高保真物理仿真，对栖息过程中关键参数（如接近速度、姿态误差、接触力变化等）进行分析，为控制策略设计提供定量依据。

在感知与规划层面，结合机载视觉与惯性测量信息，研究适用于微型无人机的栖息目标识别与相对位姿估计方法，并在此基础上设计满足姿态约束与安全距离要求的接近轨迹规划策略，使无人机能够在受限空间内稳定、可重复地完成栖息前的最终接近动作。

在控制策略设计层面，针对飞行—接触—附着过程中动力学特性突变与约束条件切换的问题，引入分层控制与混合系统控制思想，通过有限状态机或等效机制对不同工作阶段进行统一管理。在自由飞行阶段采用成熟的飞行控制方法保证轨迹跟踪性能，在最终接近与接触阶段引入更具柔顺性的控制策略，以降低接触冲击并提升附着成功率。

在机构与系统集成层面，将基于 VSERF 的柔性附着模块与无人机平台进行一体化设计，重点关注系统重量、能耗与可靠性之间的权衡，确保栖息功能的引入不会显著削弱无人机原有的飞行性能。同时，通过对附着与脱附时序的合理设计，实现对能量消耗的最小化。

最后，通过室内与半开放环境下的实验验证，对所提出方法在不同表面类型、空间构型和外部扰动条件下的性能进行系统评估，从而验证研究思路的有效性与工程可行性，并为后续实际应用奠定基础。
